% GENERATED FILE. Edit canonical lessons, then run npm run generate:books.
% canonical-source-hash: fnv1a64:1137a3732d39d285
% canonical-lessons: JA-C264-ochiru, JA-C264-butsukeru, JA-C264-houtai, JA-C264-totemo, JA-C264-shibaraku

\chapter{A Bad Fall}
\label{ch:ja-a-bad-fall}
\hlchaptermodality{264}

\begin{chapteropening}
\textbf{By the end of this chapter:} \emph{I can say how someone fell and got hurt.}

Complete the last lesson of chapter 264: I can say how someone fell and got hurt.
\end{chapteropening}

\section[\texorpdfstring{\ja{おちる} (ochiru)}{おちる (ochiru)}]{\ja{おちる} (ochiru) — to fall; to drop}

\label{lesson:JA-C264-ochiru}



\begin{quote}
Before the new one: say the Japanese for to panic, then the Japanese for at once.
\end{quote}

\subsection*{You'll want to know: \ja{おちる}}
\textbf{\ja{おちる}} — \emph{ochiru} — \textquotedblleft{}to fall; to drop\textquotedblright{}.

\begin{quote}
\textbf{\ja{こどもが} \ja{かいだんから} \ja{おちました。}} — \emph{kodomo ga kaidan kara ochimashita} — \textquotedblleft{}A child fell down the stairs.\textquotedblright{}
\end{quote}

Polite \textbf{\ja{おちます}}, \emph{ochimasu}; an ichidan verb. Past (polite): \textbf{\ja{こどもが} \ja{おちました。}}, \emph{kodomo ga ochimashita}.

\begin{cousinweb}
\textbf{\ja{おちる}} (\emph{ochiru}) is a native Japanese word (wago), inherited rather than borrowed from Chinese; its kanji was attached later as a kun'yomi (JMdict).
\end{cousinweb}

\begin{grammarlens}[title={What you've built}]
One more word for telling what happened.
\end{grammarlens}

\subsection*{Your turn}
\begin{itemize}
\raggedright
  \item \emph{Say it:} \emph{ochiru}
  \item \emph{Say it:} \emph{ochiru}, once more
  \item \emph{Say it:} \emph{kodomo ga kaidan kara ochimashita} — A child fell down the stairs.
  \item \emph{Recall:} say the Japanese for to panic, then the Japanese for at once, then say \emph{ochiru} again
\end{itemize}

\begin{tcolorbox}[breakable,colback=teal!4,colframe=teal!35!black,title={Before you move on}]
What does \ja{おちる} mean? (\textquotedblleft{}To fall; to drop\textquotedblright{}.) Say it once more.
\end{tcolorbox}
\section[\texorpdfstring{\ja{ぶつける} (butsukeru)}{ぶつける (butsukeru)}]{\ja{ぶつける} (butsukeru) — to bump (something against); to hit}

\label{lesson:JA-C264-butsukeru}



\begin{quote}
Before the new one: say the Japanese for at once, then the Japanese for to fall.
\end{quote}

\subsection*{You'll want to know: \ja{ぶつける}}
\textbf{\ja{ぶつける}} — \emph{butsukeru} — \textquotedblleft{}to bump (something against); to hit\textquotedblright{}.

\begin{quote}
\textbf{\ja{こどもが} \ja{つくえに} \ja{あたまを} \ja{ぶつけました。}} — \emph{kodomo ga tsukue ni atama o butsukemashita} — \textquotedblleft{}A child banged his head on a desk.\textquotedblright{}
\end{quote}

Polite \textbf{\ja{ぶつけます}}, \emph{butsukemasu}; an ichidan verb. Past (polite): \textbf{\ja{あにが} \ja{あたまを} \ja{ぶつけました。}}, \emph{ani ga atama o butsukemashita}.

\begin{cousinweb}
\textbf{\ja{ぶつける}} (\emph{butsukeru}) is a native Japanese word (wago), inherited rather than borrowed from Chinese, and it is usually written in kana (JMdict). It is the transitive partner of \emph{butsukaru}, \textquotedblleft{}to bump into\textquotedblright{}.
\end{cousinweb}

\begin{grammarlens}[title={What you've built}]
One more word for telling what happened.
\end{grammarlens}

\subsection*{Your turn}
\begin{itemize}
\raggedright
  \item \emph{Say it:} \emph{butsukeru}
  \item \emph{Say it:} \emph{butsukeru}, once more
  \item \emph{Say it:} \emph{kodomo ga tsukue ni atama o butsukemashita} — A child banged his head on a desk.
  \item \emph{Recall:} say the Japanese for at once, then the Japanese for to fall, then say \emph{butsukeru} again
\end{itemize}

\begin{tcolorbox}[breakable,colback=teal!4,colframe=teal!35!black,title={Before you move on}]
What does \ja{ぶつける} mean? (\textquotedblleft{}To bump (something against); to hit\textquotedblright{}.) Say it once more.
\end{tcolorbox}
\section[\texorpdfstring{\ja{ほうたい} (hōtai)}{ほうたい (hōtai)}]{\ja{ほうたい} (hōtai) — a bandage}

\label{lesson:JA-C264-houtai}



\begin{quote}
Before the new one: say the Japanese for to fall, then the Japanese for to bump (something against).
\end{quote}

\subsection*{You'll want to know: \ja{ほうたい}}
\textbf{\ja{ほうたい}} — \emph{hōtai} — \textquotedblleft{}a bandage\textquotedblright{}.

\begin{quote}
\textbf{\ja{ほうたいは} \ja{けがの} \ja{ばしょに} \ja{よります。}} — \emph{hōtai wa kega no basho ni yorimasu} — \textquotedblleft{}The kind of bandage depends on where the wound is.\textquotedblright{}
\end{quote}

\begin{cousinweb}
\textbf{\ja{ほうたい}} (\emph{hōtai}) is Sino-Japanese (kango): two Chinese characters read with their on'yomi, readings borrowed from Chinese. They are \emph{hō}, \textquotedblleft{}wrap\textquotedblright{} (Middle Chinese \emph{paew}) and \emph{tai}, \textquotedblleft{}belt\textquotedblright{} (\emph{tajH}) (KANJIDIC2; Middle Chinese in Baxter's transcription).
\end{cousinweb}

\begin{grammarlens}[title={What you've built}]
One more word for telling what happened.
\end{grammarlens}

\subsection*{Your turn}
\begin{itemize}
\raggedright
  \item \emph{Say it:} \emph{hōtai}
  \item \emph{Say it:} \emph{hōtai}, once more
  \item \emph{Say it:} \emph{hōtai wa kega no basho ni yorimasu} — The kind of bandage depends on where the wound is.
  \item \emph{Recall:} say the Japanese for to fall, then the Japanese for to bump (something against), then say \emph{hōtai} again
\end{itemize}

\begin{tcolorbox}[breakable,colback=teal!4,colframe=teal!35!black,title={Before you move on}]
What does \ja{ほうたい} mean? (\textquotedblleft{}A bandage\textquotedblright{}.) Say it once more.
\end{tcolorbox}
\section[\texorpdfstring{\ja{とても} (totemo)}{とても (totemo)}]{\ja{とても} (totemo) — very; really}

\label{lesson:JA-C264-totemo}



\begin{quote}
Before the new one: say the Japanese for to bump (something against), then the Japanese for a bandage.
\end{quote}

\subsection*{You'll want to know: \ja{とても}}
\textbf{\ja{とても}} — \emph{totemo} — \textquotedblleft{}very; really\textquotedblright{}.

\begin{quote}
\textbf{\ja{あしが} \ja{とても} \ja{いたいです。}} — \emph{ashi ga totemo itai desu} — \textquotedblleft{}My leg really hurts.\textquotedblright{}
\end{quote}

\begin{cousinweb}
\textbf{\ja{とても}} (\emph{totemo}) is a native Japanese word (wago), inherited rather than borrowed from Chinese, and it is usually written in kana (JMdict).
\end{cousinweb}

\begin{grammarlens}[title={What you've built}]
One more word for telling what happened.
\end{grammarlens}

\subsection*{Your turn}
\begin{itemize}
\raggedright
  \item \emph{Say it:} \emph{totemo}
  \item \emph{Say it:} \emph{totemo}, once more
  \item \emph{Say it:} \emph{ashi ga totemo itai desu} — My leg really hurts.
  \item \emph{Recall:} say the Japanese for to bump (something against), then the Japanese for a bandage, then say \emph{totemo} again
\end{itemize}

\begin{tcolorbox}[breakable,colback=teal!4,colframe=teal!35!black,title={Before you move on}]
What does \ja{とても} mean? (\textquotedblleft{}Very; really\textquotedblright{}.) Say it once more.
\end{tcolorbox}
\section[\texorpdfstring{\ja{しばらく} (shibaraku)}{しばらく (shibaraku)}]{\ja{しばらく} (shibaraku) — for a while; for some time}

\label{lesson:JA-C264-shibaraku}



\begin{quote}
Before the new one: say the Japanese for a bandage, then the Japanese for very.
\end{quote}

\subsection*{You'll want to know: \ja{しばらく}}
\textbf{\ja{しばらく}} — \emph{shibaraku} — \textquotedblleft{}for a while; for some time\textquotedblright{}.

\begin{quote}
\textbf{\ja{しばらく} \ja{びょういんに} \ja{いました。}} — \emph{shibaraku byōin ni imashita} — \textquotedblleft{}I was in hospital for some time.\textquotedblright{}
\end{quote}

\begin{cousinweb}
\textbf{\ja{しばらく}} (\emph{shibaraku}) is a native Japanese word (wago), inherited rather than borrowed from Chinese, and it is usually written in kana (JMdict).
\end{cousinweb}

\begin{grammarlens}[title={What you've built}]
One more word for telling what happened.
\end{grammarlens}

\subsection*{Your turn}
\begin{itemize}
\raggedright
  \item \emph{Say it:} \emph{shibaraku}
  \item \emph{Say it:} \emph{shibaraku}, once more
  \item \emph{Say it:} \emph{shibaraku byōin ni imashita} — I was in hospital for some time.
  \item \emph{Recall:} say the Japanese for a bandage, then the Japanese for very, then say \emph{shibaraku} again
\end{itemize}

\begin{tcolorbox}[breakable,colback=teal!4,colframe=teal!35!black,title={Before you move on}]
What does \ja{しばらく} mean? (\textquotedblleft{}For a while; for some time\textquotedblright{}.) Say it once more.
\end{tcolorbox}
